\section{Introduction}

In recent years there have been notable advances in the use of machine learning (ML) for drug discovery, including molecule generation and property prediction \citep{dara2022machine}.
Despite ceaseless progress in deep learning,
conventional methods such as support vector machines or random forest trained on \emph{molecular fingerprints}
are still competitive in the low-data regime \citep{walters2020applications, stanley2021fs}.
These fingerprints essentially encode fragments from a molecule into a sparse vector,
thereby compactly representing a large number of molecular substructures \citep{david2020molecular}.
They are extensively used in virtual screening for substructure and similarity searches as well as an input for ML models \citep{cereto2015molecular, granda2018controlling}.

The \emph{Tanimoto coefficient}
(also known as the \emph{Jaccard index})
stands out as a natural way to compare such fingerprints.
This coefficient is most commonly expressed as a function on sets $T_S$
or as a function of non-negative vectors $T_{MM}$
\citep{jaccard1912distribution,tanimoto1958elementary,ralaivola2005graph,costa2021further,tan2016introduction}:
\numeq{
\label{eqn:tanimoto-for-sets}
    T_S(X, X') = \frac{|X\cap X'|}{|X\cup X'|},~~ X,X'\subseteq \Omega, \qquad
    T_{MM}(x, x') = \frac{\sum_i \min(x_i, x'_i)}{\sum_i \max(x_i, x'_i)},~~ x,x'\in\mathbb{R}^d_{\geq0} \,.
}
If $a$ and $b$ are binary indicator vectors representing sets $A$ and $B$
respectively,
then $T_S(A,B)=T_{MM}(a,b)$.
Therefore $T_{MM}$ can be viewed as a generalization of $T_S$;
for this reason it is sometimes called the ``weighted Jaccard coefficient''
or min-max coefficient.

The Tanimoto coefficient is widely used in machine learning and cheminformatics
to compute similarities between molecular fingerprints
\citep{bajusz2015tanimoto,o2016comparing,miranda2021differential},
chiefly because of the following properties:
\begin{enumerate}
    \item \textbf{Clear Interpretation:}
    The value of $T_{MM}(x,x')$ represents the degree of overlap between $x$ and $x'$
    and is always between 0 and 1.
    $T(x,x')=1$ only when $x=x'$.
    \item \textbf{Kernel:}
    $T_{MM}(\cdot,\cdot)$ is positive definite \citep{gower1971general,ralaivola2005graph},
    meaning it can be used as the kernel for algorithms like support vector machines or Gaussian processes.
    \item \textbf{Metric:}
    $1-T_{MM}(x,x')$ is a valid distance metric (typically called \emph{Jaccard/Soergel distance})
    and can therefore be used in nearest-neighbour and clustering algorithms \citep{marczewski1958certain,levandowsky1971distance}.
\end{enumerate}

In this paper, we present and characterize two efficient low-rank approximations for large matrices of Tanimoto coefficients.
The first method, presented in section~\ref{sec:minmax-rf},
uses a random hash function to index a random tensor
and enjoys exceptionally low variance.
The second method, presented in section~\ref{sec:t-dp},
uses a power series expansion of the Tanimoto similarity for binary vectors.
This line of research also unexpectedly led to the discovery of a new generalization
of the Tanimoto coefficient to arbitrary vectors in $\mathbb{R}^d$, $T_{DP}$,
which is also a kernel and can be used to form a distance metric.
In section~\ref{sec:experiments} we demonstrate experimentally
that our random features are effective at approximating
Tanimoto matrices of real-world fingerprint data
and demonstrate its application to molecular property prediction and optimization problems.






